\documentclass[preview]{standalone}

\usepackage{amsmath}
\usepackage{amssymb}
\usepackage{stellar}
\usepackage{definitions}
\usepackage{bettelini}

\begin{document}

\id{power-series-boundary-behavior}
\genpage

\section{Entire functions}

\begin{snippetdefinition}{entire-function-power-series-definition}{Entire function}
    A complex \function represented by a power series with infinite radius of
    convergence is called an \emph{entire function}.
\end{snippetdefinition}

\begin{snippetexample}{entire-functions-basic-examples}{Basic entire functions}
    Every complex polynomial is entire. The exponential, sine, and cosine
    functions are also entire because their Maclaurin series have infinite
    radius of convergence.
\end{snippetexample}

\begin{snippetproposition}{entire-power-series-local-total-convergence}{Local total convergence of entire power series}
    If \(\sum_{n=0}^{\infty}a_nz^n\) has infinite radius of convergence,
    then it converges totally, and hence uniformly, on every compact subset of
    \(\complexnumbers\).
\end{snippetproposition}

\begin{snippetexample}{factorial-power-series-zero-radius-example}{A power series with zero radius}
    The power series
    \[
        \sum_{n=0}^{\infty}n!z^n
    \]
    has radius of convergence \(R=0\), as follows immediately from the ratio
    criterion.
\end{snippetexample}

\begin{snippetexample}{scaled-geometric-power-series-radius-example}{A scaled geometric power series}
    The series
    \[
        \sum_{n=0}^{\infty}\left(\frac z{27}\right)^n
    \]
    is geometric and has radius of convergence \(R=27\).
\end{snippetexample}

\begin{snippetexample}{inverse-square-power-series-closed-disk-example}{Absolute convergence on the closed disk}
    The series
    \[
        \sum_{n=1}^{\infty}\frac{z^n}{n^2}
    \]
    has radius \(R=1\) and converges absolutely at every point of the closed
    unit disk.
\end{snippetexample}

\section{Abel's theorem}

\begin{snippettheorem}{abel-power-series-boundary-theorem}{Abel's theorem for power series}
    Let \(\sum_{n=0}^{\infty}a_nz^n\) have finite positive radius of
    convergence \(R\), and suppose it converges at a boundary point
    \(z_0\), where \(|z_0|=R\). Then the series converges uniformly on the
    radial segment
    \[
        \{tz_0\suchthat 0\leq t\leq1\}.
    \]
    In the real case, if it converges at \(R\), then it converges uniformly on
    every interval \([-r,R]\) with \(0<r<R\); the analogous statement holds
    at \(-R\).
\end{snippettheorem}

\begin{snippet}{abel-power-series-boundary-theorem-interpretation}
    Abel's theorem extends uniform convergence up to a boundary point at which
    the power series converges, even when the convergence there is only
    conditional. It also permits the sum to be approached continuously along
    the corresponding radius.
\end{snippet}

\begin{snippetexample}{logarithmic-power-series-boundary-example}{Boundary behavior of the logarithmic power series}
    The series
    \[
        \sum_{n=1}^{\infty}\frac{z^n}{n}
    \]
    has radius \(R=1\). On the real boundary it converges at \(x=-1\) and
    diverges at \(x=1\). Abel's theorem gives uniform convergence on every
    interval \([-1,r]\) with \(r<1\).

    In the complex plane, the series converges at every point of the unit
    circle except \(z=1\), and the convergence is uniform on every closed arc
    that avoids \(1\).
\end{snippetexample}

\begin{snippetexample}{alternating-logarithmic-power-series-abel-example}{The alternating harmonic series from Abel's theorem}
    For \(|x|<1\),
    \[
        \sum_{n=1}^{\infty}(-1)^{n+1}\frac{x^n}{n}=\log(1+x).
    \]
    At \(x=1\) the series converges conditionally, while at \(x=-1\) it
    diverges. Abel's theorem yields uniform convergence on every interval
    \([\alpha,1]\) with \(\alpha>-1\), and therefore
    \[
        \log2
        =\lim_{x\to1^-}\log(1+x)
        =\sum_{n=1}^{\infty}\frac{(-1)^{n+1}}n.
    \]
\end{snippetexample}

\section{Limit-superior examples}

\begin{snippetexample}{power-series-sine-power-coefficients-radius-example}{Coefficients involving powers of sine}
    Consider
    \[
        \sum_{n=0}^{\infty}(\sin n)^nz^n.
    \]
    The sequence \(|\sin n|\) does not converge, but it is dense in
    \([0,1]\) after taking absolute values, so
    \[
        \limsup_{n\to\infty}\sqrt[n]{|(sin n)^n|}
        =\limsup_{n\to\infty}|\sin n|=1.
    \]
    The radius of convergence is therefore \(R=1\).
\end{snippetexample}

\begin{snippetexample}{power-series-sine-coefficients-radius-example}{Coefficients involving sine}
    For
    \[
        \sum_{n=0}^{\infty}(\sin n)z^n,
    \]
    one has
    \[
        \limsup_{n\to\infty}\sqrt[n]{|\sin n|}=1,
    \]
    and hence \(R=1\). Directly computing an ordinary limit would be less
    natural because the coefficients approach zero along some subsequences
    and stay close to one in absolute value along others.
\end{snippetexample}

\end{document}
